% Generated draft: requires mathematical and source review.
\documentclass[11pt]{article}
\usepackage[margin=25mm]{geometry}
\usepackage{fontspec}
\setmainfont{Latin Modern Roman}
\usepackage{amsmath,amssymb,mathtools,unicode-math}
\setmathfont{Latin Modern Math}
\usepackage{enumitem,needspace,xurl,hyperref}
\setlength{\parindent}{0pt}
\setlength{\parskip}{5pt}
\setlength{\emergencystretch}{3em}
\newcommand{\N}{\mathbb{N}}
\newcommand{\Z}{\mathbb{Z}}
\newcommand{\Q}{\mathbb{Q}}
\newcommand{\R}{\mathbb{R}}
\newcommand{\C}{\mathbb{C}}
\newcommand{\F}{\mathbb{F}}
\newcommand{\A}{\mathbb{A}}
\newcommand{\PP}{\mathbb{P}}
\newcommand{\E}{\mathbb{E}}
\newcommand{\eps}{\varepsilon}
\newcommand{\dd}{\,\mathrm{d}}
\begin{document}
\section*{UnsolvedMath ID: 30000593}
\textbf{Source problem code:} catalogue entry
\subsection{Definitions and mathematical statement}
Let \(M\) be a closed smooth spin manifold of dimension \(n\), and put \(\Gamma=\pi_1(M)\) in the connected convention. The first question is the proposed implication
\[
 M\text{ has infinite stable }K\text{-area}
 \quad\Longrightarrow\quad\alpha_{\mathrm{max}}(M)\ne0.
\]
The second asks for an analogous nonvanishing statement for \(\alpha_{\mathrm{red}}(M)\), at least in specified special cases such as area-enlargeability. These are retained as separate maximal and reduced questions.

In the usual Rosenberg-index setup, \(\alpha_{\mathrm{max}}(M)\) is the index of the spin Dirac operator twisted by the Mishchenko bundle with fiber the maximal group \(C^*\)-algebra of \(\Gamma\); the reduced index uses the reduced group \(C^*\)-algebra. The natural quotient
\[
 q:C^*_{\mathrm{max}}\Gamma\longrightarrow C^*_{\mathrm{red}}\Gamma
\]
induces \(q_*\alpha_{\mathrm{max}}(M)=\alpha_{\mathrm{red}}(M)\) in the compatible index theory. This does not imply that a nonzero maximal class has nonzero reduced image. The source does not specify real \(KO_n\) versus complex \(K_n\) coefficients, which can affect nonvanishing; neither is silently substituted for the other.

Infinite stable \(K\)-area is not defined in the record. Usual \(K\)-area notions detect nontrivial characteristic data with bundles of arbitrarily small curvature, while ``stable'' requires a prescribed stabilization and detection convention. The allowed bundles, curvature norm, characteristic pairing, and stabilization are missing, so a particular torus product or finite-rank versus Hilbert-module definition is not imposed.

For orientation, a common area-enlargeability condition asks, for every \(\varepsilon>0\), for a suitable covering \(\widetilde M_\varepsilon\to M\) and a nonzero-degree map to \(S^n\) whose differential contracts two-dimensional areas, \(\|\Lambda^2df\|\leq\varepsilon\), with an appropriate constant-outside-compact-set condition for a noncompact cover. The precise allowed covers, target/degree conventions and dimension hypotheses must come from the source. This is not replaced by contraction of lengths, which is the stronger ordinary enlargeability-type condition. Nor is the relation of area-enlargeability to the missing stable-\(K\)-area definition assumed without its hypotheses.

The literature assessment checked 2026-08-21 reports maximal-index nonvanishing for enlargeable spin manifolds and related positive-scalar-curvature obstructions, while leaving the full stable-\(K\)-area implications and reduced-index statement unresolved. The draft keeps the reported enlargeability result distinct from the specifically mentioned area-enlargeability question. Its partially-solved label and notation-only verification are dated metadata, not independently verified current results.
\subsection{Short English statement}
Does infinite stable \(K\)-area force a nonzero maximal Rosenberg index, and can reduced-index nonvanishing be proved under that or relevant special hypotheses such as area-enlargeability? The two completions are distinct: the quotient map can kill a nonzero class. Stable \(K\)-area, coefficient theory and covering conventions are missing.
\subsection{Sources}
\begin{itemize}
\item Local UnsolvedMath registry, permanent ID 30000593, source code catalogue entry, ``Stable K-Area and Nonvanishing Rosenberg Indices,'' including both maximal/reduced questions and the area-enlargeability special-case request.
\item \emph{Topologie (2007)}, Oberwolfach Reports, pp.\ 2579--2630, DOI 10.4171/owr/2006/43: \url{https://doi.org/10.4171/owr/2006/43}. The citation uses 2007, while the DOI and proposed-year metadata use 2006.
\item Supplied reference ``Hanke--Schick, enlargeability and index theory'': \url{https://www.uni-math.gwdg.de/schick/publ/Enlargeability_and_index_theory_final_for_JDG.pdf}. The URL is cited as supplied; no PDF has been opened.
\item Literature triage checked 2026-08-21; record updated 2026-09-16. No external source has been visited or independently checked here.
\end{itemize}
\end{document}
